% ---------------------------------------------------------------------------
% Topic1ExplanationsPart1 — pages 10–16
% Topic SC-1 Questions 1–2; Topic 1-3 Questions 1–4
% ---------------------------------------------------------------------------

\lesson{Lesson SC-1: Compare Rational Numbers}

\question{Question 1}

\blockhead{Question}
Your pet cat weighs 10.4 lbs. Your friend's rabbit weighs 10.6 lbs. Write two different statements using the less than ($<$) and greater than ($>$) symbols that compare the weights.

\checked{$10.4 < 10.6$ and $10.6 > 10.4$}

\blockhead{What the question is asking}
This question asks for two true comparison sentences relating 10.4 and 10.6 using the less-than mark ($<$) and the greater-than mark ($>$).

\blockhead{Important words}
\begin{words}
  \item A \term{decimal}has digits after the decimal point. Here both weights have one digit after the point.
  \item The \term{ones place}is the place just left of the decimal point. It counts whole ones.
  \item \term{Tenths}is the first digit place to the right of the decimal point.
  \item \term{Hundredths}is the second digit place to the right of the decimal point.
  \item \term{Place value}is the value of a digit's position (ones, tenths, hundredths, $\ldots$).
  \item The symbol $<$ means \term{less than}(smaller).
  \item The symbol $>$ means \term{greater than}(larger).
  \item To \term{compare}means to decide which is smaller or larger.
  \item To \term{flip}a comparison means write the same fact the other way ($a < b$ becomes $b > a$).
\end{words}

\blockhead{Work the problem one tiny step at a time}
\begin{steps}
  \item Line up the decimals by place value: both start with 10.
  \item Compare the tenths place: 4 tenths versus 6 tenths.
  \item Because $4 < 6$, we have $10.4 < 10.6$.
  \item Flip the comparison the other way: $10.6 > 10.4$.
  \item Both statements are true and use different symbols as required.
\end{steps}

\finalans{$10.4 < 10.6$ and $10.6 > 10.4$ (pounds OK)}

\blockhead{Why this makes sense}
The cat is lighter. Saying the smaller number is less than the larger number, and the larger is greater than the smaller, are two ways to say the same fact.

\question{Question 2}

\begin{samepage}
\blockhead{Question}
Compare the fraction $\frac{3}{4}$ with the decimal 0.75. Select all symbols that appropriately relate the two numbers.
\nopagebreak
\begin{center}
$<$\qquad $>$\qquad $=$\qquad $\neq$
\end{center}
\end{samepage}

\checked{$=$ (only)}

\blockhead{What the question is asking}
This question asks which comparison symbols are true for $\frac{3}{4}$ and 0.75.

\blockhead{Important words}
\begin{words}
  \item A \term{fraction}$\frac{3}{4}$ means 3 divided by 4.
  \item A \term{decimal}is a number written with a decimal point to show parts smaller than one whole. Here the decimal is 0.75.
  \item \term{Long division}is a written way to divide one digit at a time.
  \item A \term{remainder}is the amount left after one division step.
  \item To \term{bring down}a digit means move the next digit beside the remainder so division can continue.
  \item The symbol $=$ means \term{equal to.}
  \item The symbol $<$ means \term{less than}(strictly smaller, never equal).
  \item The symbol $>$ means \term{greater than}(strictly larger, never equal).
  \item The symbol $\neq$ means \term{not equal to.}
  \item \term{Select all}means check every symbol on the list and keep only the ones that make a
        true statement. Here the list offers just four symbols: $<$, $>$, $=$, and $\neq$.
\end{words}

\blockhead{Work the problem one tiny step at a time}
\begin{steps}
  \item Change the fraction to a decimal: $\frac{3}{4} = 3 \div 4$.
  \item Write 3 as 3.00 so the division can continue past the decimal point.
  \item Bring down the first 0. This makes 30.
  \item The number 4 fits into 30 seven times.
  \item Multiply to check that digit: $4 \times 7 = 28$.
  \item Subtract: $30 - 28 = 2$. The remainder is 2.
  \item Bring down the next 0. This makes 20.
  \item The number 4 fits into 20 five times because $4 \times 5 = 20$.
  \item Subtract: $20 - 20 = 0$. The remainder is 0.
  \item The decimal digits are 7 and 5. So $3 \div 4 = 0.75$.
  \item Compare the two numbers: $\frac{3}{4} = 0.75$ exactly. So $=$ is true.
  \item Now test each of the other three symbols offered.
  \item Test $<$: it would say $\frac{3}{4}$ is strictly smaller than 0.75. The numbers are equal,
        so $<$ is false.
  \item Test $>$: it would say $\frac{3}{4}$ is strictly larger than 0.75. The numbers are equal,
        so $>$ is false.
  \item Test $\neq$: it would say the two numbers are not equal. They are equal, so $\neq$ is false.
  \item Only one symbol out of the four makes a true statement, so select $=$ and nothing else.
\end{steps}

\finalans{$=$ (only)}

\blockhead{Why this makes sense}
The fraction $\frac{3}{4}$ and the decimal 0.75 are two ways of writing the very same number, so the only relationship that holds between them is equality. Each of the other three choices claims the numbers differ in some way --- smaller, larger, or unequal --- and every one of those claims is false when the numbers are the same.

\lesson{Lesson 1-3: Add Integers}

\question{Question 1}

\blockhead{Question}
In round 1 of a game, a player scores 10 points, then loses 4 points, then loses 8 points.

\par\noindent\runin{Part A:}Explain how to show the change on a number line.
\par\noindent\runin{Part B:}What is the final score?

\checked{Part A: start at 0, move right 10, left 4, left 8; Part B: $-2$}

\blockhead{What the question is asking}
This question asks you to show the score changes on a straight line of numbers and find the final score after $+10$, then $-4$, then $-8$.

\blockhead{Important words}
\begin{words}
  \item A \term{decimal}uses a decimal point to show parts of a whole.
  \item An \term{integer}is one of the numbers $\ldots, -2, -1, 0, 1, 2, \ldots$; it has no decimal part.
  \item A \term{positive number}is greater than zero. A \term{negative number}is less than zero.
  \item To \term{add}means join an amount to the amount you have.
  \item A \term{number line}is a straight line with numbers in order. Moving right adds a positive number. Moving left adds a negative number.
  \item A \term{unit}on the number line is one tick-mark step from one integer to the next; on this number line one unit is one point.
  \item A \term{score}of 10 points is $+10$. Losing 4 points is $-4$. Losing 8 points is $-8$.
  \item The \term{final score}is what you get after all the moves.
\end{words}

\blockhead{Work the problem one tiny step at a time}
\begin{steps}
  \item \runin{Part A:}Start at 0 on a number line.
  \item Move right 10 units to land on 10 (score 10 points).
  \item Move left 4 units to land on 6 (lose 4).
  \item Move left 8 units to land on $-2$ (lose 8).
  \item \runin{Part B:}Add the integers: $10 + (-4) + (-8)$.
  \item First: $10 + (-4) = 6$.
  \item Then: $6 + (-8) = -2$.
\end{steps}

\begin{center}
\begin{tikzpicture}[x=0.62cm,y=0.62cm]
  % Number line ticked at every integer, so one tick-mark step is exactly one unit.
  \draw[{Stealth[length=2.4mm]}-{Stealth[length=2.4mm]}] (-4.7,0) -- (12.7,0);
  \foreach \x in {-4,-3,...,12}{%
    \draw (\x,0.16) -- (\x,-0.16);
    \node[below=1pt,font=\footnotesize] at (\x,-0.16) {$\x$};}
  % The four points the player stands on: start 0, then 10, then 6, then -2.
  \foreach \p in {0,10,6,-2}{\fill (\p,0) circle (2.4pt);}
  % Guides carrying each landing up to the move that starts there.
  \draw[gray,dashed] (0,0) -- (0,1.0);
  \draw[gray,dashed] (10,0) -- (10,2.3);
  \draw[gray,dashed] (6,0) -- (6,3.6);
  \draw[gray,dashed] (-2,0) -- (-2,3.6);
  % The three moves of Part A.
  \draw[-{Stealth[length=2.4mm]},thick] (0,1.0) --
    node[above,font=\footnotesize] {1. right 10 (score 10)} (10,1.0);
  \draw[-{Stealth[length=2.4mm]},thick] (10,2.3) --
    node[above,font=\footnotesize] {2. left 4 (lose 4)} (6,2.3);
  \draw[-{Stealth[length=2.4mm]},thick] (6,3.6) --
    node[above,font=\footnotesize] {3. left 8 (lose 8)} (-2,3.6);
\end{tikzpicture}
\end{center}

\finalans{Part A: start at 0, move right 10, left 4, left 8; Part B: $-2$}

\blockhead{Why this makes sense}
Gains go right; losses go left. Ending left of zero means a negative score. Check: $10 - 4 - 8 = -2$.

\question{Question 2}

\blockhead{Question}
A student starts a game with 15 points. They lose 20 points in the first round and then gain 3 points in the second round. What is the student's final score?

\checked{$-2$}

\blockhead{What the question is asking}
This question asks for the final score after starting at 15, losing 20, then gaining 3.

\blockhead{Important words}
\begin{words}
  \item An \term{operation}is a math action. Add joins amounts. Subtract takes away or finds a difference. Multiply makes equal groups. Divide splits into equal groups.
  \item An \term{expression}is a math phrase made of numbers and operation signs.
  \item A \term{decimal}uses a decimal point to show parts of a whole.
  \item An \term{integer}is one of the numbers $\ldots, -2, -1, 0, 1, 2, \ldots$; it has no decimal part.
  \item To \term{lose}points means add a \term{negative}integer.
  \item To \term{gain}points means add a \term{positive}integer.
  \item On a \term{number line,}moving left adds a negative; moving right adds a positive.
  \item The \term{final score}is the value after all the additions are done.
  \item An \term{addend}is a number being added. Here the addends are 15, $-20$, and 3.
\end{words}

\blockhead{Work the problem one tiny step at a time}
\begin{steps}
  \item Write the expression: $15 + (-20) + 3$.
  \item First combine $15 + (-20)$. Start at 15. Move left 20.
  \item From 15, move left 15 units to reach 0. That uses 15 of the 20 leftward steps.
  \item Steps still left: $20 - 15 = 5$. From 0, move left 5 more to land on $-5$.
  \item So $15 + (-20) = -5$.
  \item Now add 3: $-5 + 3$.
  \item Start at $-5$. Move right 3. Land on $-2$.
  \item So the final score is $-2$.
  \item Check: $15 - 20 = -5$, then $-5 + 3 = -2$. Same answer.
\end{steps}

\finalans{$-2$}

\blockhead{Why this makes sense}
Losing more than you started with puts you below zero, then gaining 3 only climbs back a little: $-5 + 3 = -2$. Check: $15 - 20 + 3 = -2$.

\question{Question 3}

\blockhead{Question}
A bank account starts with $\$500$. The bank deducts a monthly maintenance fee of $\$12.50$. If there are no other deposits or withdrawals, how much money is in the account after 6 months?

\checked{$\$425$}

\blockhead{What the question is asking}
This question asks how much money is left after six monthly charges of $\$12.50$.

\blockhead{Important words}
\begin{words}
  \item A \term{balance}is how much money is in the account.
  \item A \term{negative}number is less than zero.
  \item To \term{multiply}means make equal groups. For example, $6 \times \$12.50$ means 6 equal groups of $\$12.50$.
  \item To \term{subtract}means take away or find what remains.
  \item A \term{fee}is money taken away. Taking away means subtract (or add a negative).
  \item A \term{deposit}adds money. A \term{withdrawal}removes money. Here there are none besides the fee.
  \item \term{Monthly}means once each month.
\end{words}

\blockhead{Work the problem one tiny step at a time}
\begin{steps}
  \item Fee each month: $\$12.50$.
  \item Number of months: 6.
  \item Total fees: multiply $6 \times \$12.50$.
  \item Compute $6 \times \$12 = \$72$ and $6 \times \$0.50 = \$3$, then $\$72 + \$3 = \$75$. Total fees $= \$75$.
  \item Subtract from the start: $\$500 - \$75$.
  \item Compute $\$500 - \$75 = \$425$.
\end{steps}

\finalans{$\$425$}

\blockhead{Why this makes sense}
Six fees of $\$12.50$ make $\$75$. Starting at $\$500$ and removing $\$75$ leaves $\$425$. A common mistake is subtracting $\$12.50$ only once.

\question{Question 4}

\blockhead{Question}
Evaluate: $-12 + 19 =$

\checked{7}

\blockhead{What the question is asking}
This question asks for the value of $-12 + 19$. You start left of zero and move right by a larger amount.

\blockhead{Important words}
\begin{words}
  \item An \term{operation}is a math action. Add joins amounts. Subtract takes away or finds a difference. Multiply makes equal groups. Divide splits into equal groups.
  \item An \term{expression}is a math phrase made of numbers and operation signs.
  \item To \term{evaluate}means find the value of an expression.
  \item A \term{unit}is one single step or amount.
  \item A \term{number line}puts numbers in order: negatives left of zero, positives right of zero.
  \item Adding a \term{positive}number on a number line means move \textbf{right}.
  \item $-12$ is 12 units \textbf{left} of zero. $+19$ means move 19 units right from wherever you start.
  \item \term{Absolute value}means distance from zero. Here $\abs{-12} = 12$ and $\abs{19} = 19$.
  \item An \term{addend}is a number being added. In $-12 + 19$, the addends are $-12$ and 19.
  \item The \term{sum}is the answer to an addition problem.
  \item When the positive addend is larger than the absolute value of the negative start, the sum is positive.
\end{words}

\blockhead{Work the problem one tiny step at a time}
\begin{steps}
  \item Start at $-12$ on the number line (12 left of zero).
  \item Add 19 means move right 19 units from $-12$.
  \item First, move right 12 units to reach 0. That uses up 12 of the 19 steps.
  \item Steps used so far: 12. Steps still left: $19 - 12 = 7$.
  \item From 0, move right the remaining 7 units.
  \item Land on $+7$. So $-12 + 19 = 7$.
  \item Check another way: the larger absolute value is 19 (positive), and $19 - 12 = 7$, so the sum is $+7$.
\end{steps}

\finalans{7}

\blockhead{Why this makes sense}
You had to ``pay back'' 12 just to get back to zero, then 7 steps remained to the right. Check: $19 - 12 = 7$, and the answer is positive because 19 is larger than 12.
